\documentclass[11pt]{letter}
\usepackage[utf8]{inputenc}
\usepackage{geometry}
\geometry{letterpaper, margin=1in}

\address{Rafael D. De Paz \\ Independent Researcher \\ me@rdepaz.com}
\signature{Rafael D. De Paz}

\begin{document}
\begin{letter}{Editorial Board \\ Journal of Mathematical Physics \\ American Institute of Physics}

\opening{Dear Editors,}

I am submitting for your consideration the manuscript entitled \emph{``The Discrete Frequency Divider: Connecting Proton Compton Oscillation to the 136.1 Hz Carrier Wave via the Icosahedral 60R-Grid.''}

This work establishes a novel mathematical connection between the proton Compton frequency ($f_{proton} \approx 2.27 \times 10^{23}$ Hz) and the macroscopic 136.1 Hz resonance through the discrete geometric structure of the Poincar\'{e} homology sphere $S^3/2I$. By modeling the binary icosahedral group as a cascading frequency divider with 12 stages, we demonstrate that 136.1 Hz emerges as a necessary geometric echo of the proton's quantum oscillation --- not a cultural artifact but a topological inevitability.

The paper bridges spectral geometry, harmonic analysis, and the Hilbert-P\'{o}lya conjecture in a manner that has not previously appeared in the literature. All predictions are falsifiable via cosmological topology measurements (Planck/future CMB missions) and proton charge radius refinements.

This manuscript has not been submitted elsewhere and all work is original.

\closing{Respectfully,}

\end{letter}
\end{document}
